\begin{theorem}[G\"odel--Gentzen negative translation]
\label{thm:negative-translation}
Assign to every first-order formula \(A\) a formula \(A^N\) by recursion:
\(\bottom^N\defeq\bottom\), and this clause takes priority in the sense
that \(\bottom\) is formed by its own grammar clause, is not an atomic
formula (\cref{sec:fol-syntax-semantics-02}), and is not treated by the next
clause; \(A^N\defeq\neg\neg A\) for atomic \(A\); the translation commutes
with \(\land\), \(\to\), and
\(\forall\); \((A\lor B)^N\defeq\neg(\neg A^N\land\neg B^N)\); and
\((\exists x\,A)^N\defeq\neg\forall x\,\neg A^N\).  Here \(\neg\) is the
abbreviation \(\neg A\defeq A\to\bottom\) of \cref{sec:lk-lj-02}, so a
negated formula is treated by the \(\to\) clause; by unique readability
(\cref{thm:fol-unique-readability}), exactly one clause applies to each
formula, and the recursion is well defined.  Then
\(\LK\derives\Gamma\Rightarrow A\) iff
\(\LJ\derives\Gamma^N\Rightarrow A^N\), where \(\Gamma^N\) translates
every member.
\textit{Status: cited theorem.}
\end{theorem}

See \cite{TroelstraVanDalen1988}.

